\documentclass[12pt]{article}

% ------------------------------------------------------------------------------
% Packages
% ------------------------------------------------------------------------------
\usepackage[margin=1in]{geometry}
\usepackage{amsmath, amssymb, amsthm}
\usepackage{natbib}
\usepackage{hyperref}
\usepackage{booktabs}
\usepackage{graphicx}
\usepackage{setspace}
\usepackage{microtype}

\doublespacing

% ------------------------------------------------------------------------------
% Title
% ------------------------------------------------------------------------------
\title{Demand Shocks in the Italian Government Bond Market}

\author{Alessandro Calzolaio%
  \thanks{University of Trento.
    Email: \texttt{alessandro.calzolaio@unitn.it}.}
}

\date{This version: \today}

% ------------------------------------------------------------------------------
\begin{document}
% ------------------------------------------------------------------------------

\maketitle

\begin{abstract}
[Abstract placeholder. Identification strategy, main findings, and contribution
to go here.]
\end{abstract}

\vspace{1em}
\noindent\textbf{JEL Codes:} G12, G14, E43, E58 \\
\noindent\textbf{Keywords:} government bonds, auction demand shocks,
spillovers, Italian Treasury, BTP

\clearpage

% ------------------------------------------------------------------------------
\section{Introduction}
\label{sec:intro}
% ------------------------------------------------------------------------------

[Introduction placeholder.]

% Literature positioning, contribution, and roadmap to go here.

% ------------------------------------------------------------------------------
\section{Institutional Background and Data}
\label{sec:data}
% ------------------------------------------------------------------------------

\subsection{Italian Treasury Auctions}

[Description of BTP auction mechanism, timing, and Bank of Italy role.]

\subsection{Data Sources}

[BTC\_BoI series, comovement data (COMOVEMENT.xlsx), tick cache, sample period.]

% ------------------------------------------------------------------------------
\section{Identification of Demand Shocks}
\label{sec:identification}
% ------------------------------------------------------------------------------

\subsection{Shock Definition}

Following \citet{Lengyel2021_demand_shocks}, we measure the demand shock as the
log price change in the narrow 10:59--11:15 window surrounding the announcement
of Italian government bond auction results.

\[
  \varepsilon_t^{(m)} = \ln P_{11:15,t}^{(m)} - \ln P_{10:59,t}^{(m)}
\]

where $m \in \{3\text{Y}, 5\text{Y}, 10\text{Y}, 15\text{Y}, 30\text{Y}\}$
denotes maturity and $t$ indexes auction dates.

\subsection{Normalization and Pooling}

[Normalization by maturity-specific volatility; pooled long panel construction.]

% ------------------------------------------------------------------------------
\section{Exogeneity Tests}
\label{sec:exogeneity}
% ------------------------------------------------------------------------------

[Panel C regressions: demand shock on fundamentals/policy expectations.
Orthogonality with ECB surprises (EA-MPD, \citealt{Altavilla2019_eampd}).]

% ------------------------------------------------------------------------------
\section{Spillover Results}
\label{sec:spillovers}
% ------------------------------------------------------------------------------

\subsection{Cross-Market Comovement}

[Panels A \& B: OLS-HAC and Huber-T estimates for corporate/private debt
and equities.]

\subsection{Magnitude Analysis}

[Rescaled effects in log-price basis-point units.]

% ------------------------------------------------------------------------------
\section{Robustness}
\label{sec:robustness}
% ------------------------------------------------------------------------------

[Symmetric window sensitivity; alternative HAC lags; winsorization; day-before
placebo; ECB meeting-day exclusion; crisis + COVID exclusion; Wald pooling test.]

% ------------------------------------------------------------------------------
\section{Conclusion}
\label{sec:conclusion}
% ------------------------------------------------------------------------------

[Conclusion placeholder.]

% ------------------------------------------------------------------------------
\bibliographystyle{aer}
\bibliography{../Bibliography_base}
% ------------------------------------------------------------------------------

\clearpage
\appendix

\section{Additional Tables and Figures}
\label{app:tables}

[Appendix placeholder.]

% ------------------------------------------------------------------------------
\end{document}
% ------------------------------------------------------------------------------
